\subsection{Elementare Zeilenumformungen}
\label{subsec:linear_systems_elementary_row_operations}

Die folgenden Umformungen ändern die Lösungsmenge nicht:
\[
  \begin{alignedat}{2}
    Z_i&\leftrightarrow Z_j
      &\qquad&\text{zwei Zeilen vertauschen},\\[0.45em]
    Z_i&\leftarrow \lambda Z_i,\quad \lambda\neq0
      &&\text{eine Zeile skalieren},\\[0.45em]
    Z_i&\leftarrow Z_i+\lambda Z_j
      &&\text{ein Vielfaches einer anderen Zeile addieren}.
  \end{alignedat}
\]

Sie heißen \emph{elementare Zeilenumformungen}. Jede von ihnen ist umkehrbar;
deshalb entsteht jeweils ein äquivalentes LGS mit derselben Lösungsmenge.

\begin{ExampleBox}{Eine Variable eliminieren}{linear_systems_eliminate_variable}
  \begin{center}
    \begin{tikzpicture}[
      >=Latex,
      gaussmatrix/.style={
        matrix of math nodes,
        left delimiter={[},right delimiter={]},
        nodes={minimum width=0.75cm,minimum height=0.58cm,anchor=center},
        column sep=1pt,row sep=1pt
      },
      operation/.style={font=\normalsize,anchor=west}
    ]
      \matrix (M0) [gaussmatrix] {
        1&1&2\\
        2&3&5\\
      };
      \draw[line width=0.6pt]
        ($(M0-1-2.north east)!0.5!(M0-1-3.north west)+(0,0.12)$)--
        ($(M0-2-2.south east)!0.5!(M0-2-3.south west)-(0,0.12)$);
      \node[operation] (op) at ($(M0.east)+(1.45,0)$ |- M0-2-3.east)
        {$Z_2\leftarrow Z_2-2Z_1$};
      \draw[->,line width=0.7pt]
        ($(M0.east)+(0.15,0)$ |- M0-2-3.east)--($(op.west)+(-0.15,0)$);

      \matrix (M1) [gaussmatrix,below=0.8cm of M0] {
        1&1&2\\
        0&1&1\\
      };
      \draw[line width=0.6pt]
        ($(M1-1-2.north east)!0.5!(M1-1-3.north west)+(0,0.12)$)--
        ($(M1-2-2.south east)!0.5!(M1-2-3.south west)-(0,0.12)$);
    \end{tikzpicture}
  \end{center}
  Die neue zweite Zeile enthält $x_1$ nicht mehr. Aus $x_2=1$ folgt durch
  Rückwärtseinsetzen $x_1=1$.
\end{ExampleBox}
